%% CS6886 course project report -- two-page limit excluding references.
%% Build:  tools/tectonic/tectonic.exe submission/report.tex
\documentclass[10pt,twocolumn]{article}

\usepackage[margin=0.68in]{geometry}
\usepackage{times}
\usepackage[T1]{fontenc}
\usepackage{amsmath}
\usepackage{booktabs}
\usepackage{graphicx}
\usepackage{microtype}
\usepackage[hidelinks]{hyperref}
\usepackage{titlesec}
\usepackage{enumitem}
\usepackage{xcolor}

\definecolor{navy}{HTML}{1E2761}

% compress vertical space so the content fits two pages
\titlespacing*{\section}{0pt}{5pt}{2pt}
\titlespacing*{\subsection}{0pt}{4pt}{1pt}
\setlength{\parskip}{2pt}
\setlength{\abovedisplayskip}{3pt}
\setlength{\belowdisplayskip}{3pt}
\setlist[itemize]{topsep=1pt,itemsep=0pt,leftmargin=12pt,parsep=0pt}
\renewcommand{\baselinestretch}{0.97}

\titleformat{\section}{\normalfont\bfseries\color{navy}\large}{\thesection}{0.5em}{}
\titleformat{\subsection}{\normalfont\bfseries\color{navy}\normalsize}{\thesubsection}{0.5em}{}

\newcommand{\best}[1]{\textbf{#1}}

\begin{document}

\twocolumn[
\begin{@twocolumnfalse}
\centering
{\LARGE\bfseries\color{navy} Which Bits Can We Drop?\par}
\vspace{3pt}
{\large Component-Wise Quantization, Mixed-Precision Selection and QAT for a Learned Image Codec\par}
\vspace{4pt}
{\normalsize Harshith (and team) \quad\textbar\quad CS6886 Final Project \quad\textbar\quad November 2026\par}
\vspace{2pt}
{\small Code: \url{https://github.com/harshith0518/CS6886-PULSE-Quantization}\par}
\vspace{8pt}
\begin{minipage}{0.92\textwidth}
\small
\textbf{Abstract.}
A learned image codec stores one weight file that must serve eight quality settings, so the usual
recipe---pick one precision for the whole model---ignores that the model is wildly unbalanced.
We measure how much each of nineteen components tolerates seven formats across eight quality
settings and two datasets (69{,}233 measurements), and feed those costs to a mixed-integer program
that chooses one precision per component. Result: \best{21.677 MB} against 54.688 MB FP32, a
\best{60.4\%} reduction at a \best{0.94\%} worst-case rate-equivalent penalty and $+0.48$ to
$+0.55\%$ BD-rate on three datasets. Against uniform INT8---which yields a \emph{smaller} file at
26.18\% penalty---the mixed assignment is \best{39$\times$} better quality at comparable size.
We then test quantization-aware fine-tuning (QAT): it removes \best{61--98\%} of the damage on
hand-chosen assignments and \best{exactly 0\%} on the solver-optimised one, so QAT repairs what an
assignment leaves on the table, and an optimal assignment leaves nothing.
\end{minipage}
\vspace{8pt}
\end{@twocolumnfalse}
]

\section{Introduction}
Post-training quantization is the standard way to shrink a neural codec, and the standard practice
is a single precision for the whole model because per-layer analysis costs time. For
\emph{pulse-mse}~\cite{pulse} that reflex is expensive: 85.6\% of its 14.18M parameters sit in six
encoder blocks, while the entropy model---which dominates decode latency---holds 9.4\%, and the
708k-parameter decoder is only 5.0\% of the network yet contains a 4.7k block that loses 8\% at 8 bits.
Sensitivity is therefore not monotone in size: a 2.89M-parameter encoder block tolerates 8-bit weights
at 0.5\% penalty while that 4.7k renderer MLP loses 8\%.

We therefore ask three questions. \textbf{(Q1)} Which components tolerate which precision, at which
quality setting? \textbf{(Q2)} Given those measurements, what is the smallest per-component
assignment that stays inside a quality budget? \textbf{(Q3)} Can QAT recover the residual loss, and
where does it stop paying? A fourth, unplanned question proved the most useful: why does
quantization not make this codec faster?

\section{Prior work}
Integer quantization with per-output-channel symmetric scales and straight-through gradient
estimation is established practice~\cite{jacob,bengio}. Learned image compression is represented here
by PULSE~\cite{pulse}. Mixed-precision allocation as an integer program is a standard
formulation~\cite{highs}; we solve it with HiGHS through SciPy. For rate-distortion comparison we
follow Bj{\o}ntegaard delta-rate~\cite{bjontegaard}, deviating only in using monotone piecewise-cubic
interpolation~\cite{pchip} because we have eight points per curve rather than the four the original
method assumes.

\section{Idea and implementation}
\subsection{Measuring sensitivity}
We quantize \emph{one} component at a time, leaving every other component and all activations in
FP32, and run the full codec over every image at every quality setting. Seven formats, nineteen
components, two datasets (Kodak~\cite{kodak} screening, CLIC Professional~\cite{clic} confirmation)
and eight settings give 69{,}233 rows with zero failures and weight restoration verified by hash. An
FP32 gate reproduces the reference bitstream and reconstruction hashes exactly at all eight settings,
so any difference is attributable to the quantized component.

Quality is summarised by a rate-equivalent penalty that folds a quality change and a bitrate change
into one number,
\begin{equation}
\text{penalty}=100\left[\left(1+\tfrac{\Delta\text{bpp}}{100}\right)2^{-\Delta\text{PSNR}/\sigma}-1\right],
\end{equation}
where $\sigma$ is the local slope of the reference curve. Zero means indistinguishable from FP32.
The optimiser constrains the worst setting: the maximum over the eight settings of the mean over
images.

Optimising over the sum of single-component costs assumes additivity, so we tested it: a uniform INT8
model's measured penalty against the sum of its nineteen component penalties gives a ratio of
0.96--1.05 on Kodak and 1.02--1.11 on CLIC. Additivity holds. A trap worth recording: summing each
component's \emph{own} worst setting inflates the total to 65.3\% where the true value at a single
setting is 26.7\%; costs must be summed per setting, then maximised.

\subsection{The allocation problem}
Because one weight file serves all eight settings, a per-setting choice would not be deployable. The
model has one binary variable per (component, format) pair, shared across settings, with one quality
constraint per setting:
\begin{equation*}
\begin{aligned}
\min\;&\textstyle\sum_{c,f}\text{size}(c,f)\,x_{c,f}\\
\text{s.t.}\;&\textstyle\sum_{c,f}\text{pen}(c,f,q)\,x_{c,f}\le \epsilon \quad \forall q\in\{0,\dots,7\}\\
&\textstyle\sum_{f} x_{c,f}=1 \quad \forall c .
\end{aligned}
\end{equation*}
The objective is bytes; deliberately not latency, because no per-format latency table exists for this
GPU and an objective using one would optimise invented data. A greedy baseline is solved alongside.

\subsection{Quantization-aware fine-tuning}
QAT here means straight-through quantized convolutions followed by minimising distortion subject to
holding the bitrate at its pre-quantization value. A fixed rate penalty fails in both directions---too
weak and the model drifts down the RD curve, too strong and it destroys quality to hit the target---so
we use an augmented Lagrangian whose multiplier is updated from the measured rate error on a fixed
validation set:
\begin{equation*}
\mathcal{L}=\frac{D}{D_{\text{ref}}}+\mu\left(\frac{R}{R_{\text{ref}}}-1\right),\qquad
\mu\leftarrow\mathrm{clip}\!\left(\mu+\alpha\,e_{\text{rate}}\right).
\end{equation*}
Entropy components are \emph{frozen}: the shipped bundle holds precomputed integer CDF tables that
would no longer describe the weights, breaking the encoder/decoder contract. The saved checkpoint must
beat its own starting point, so a bad run cannot ship a worse model. Training uses 100 DIV2K~\cite{div2k}
images disjoint from every test set, 12{,}000 steps at 512-pixel crops, about 100 minutes per tier on
one 4\,GB laptop GPU---2.2\% of the optimiser steps and 0.2\% of the sample budget of
pretraining~\cite{pulse}. This is a repair operation, not retraining.

\section{Evaluation}
\subsection{Against the baseline}
Table~\ref{tab:frontier} reports all three operating points, measured on all three test sets at all
eight settings against the same FP32 reference. The solver predicted 0.9715\% for the champion and
measured 0.94\%: the analytic model errs conservative on every dataset. Set14 was never used for
fitting and holds, which is the best evidence that the assignment transfers.

\begin{table}[t]
\centering
\footnotesize
\setlength{\tabcolsep}{3pt}
\caption{Measured operating points. Worst-setting penalty (\%), lower is better. FP32 is 54.688 MB;
BD-rate for the champion is $+0.48/+0.55/+0.49\%$.}
\label{tab:frontier}
\begin{tabular}{lrrrr}
\toprule
Point & MB & Save & Kodak & CLIC \\
\midrule
1\% champion      & \best{21.677} & $-60.4\%$ & \best{0.67} & \best{0.85} \\
1.8\% variant     & 18.664 & $-65.9\%$ & 1.27 & 1.60 \\
INT8 tier + QAT   & \best{14.401} & $-73.7\%$ & 1.85 & 2.87 \\
\midrule
uniform INT8      & 13.904 & $-74.6\%$ & 26.18 & 56.70 \\
\bottomrule
\end{tabular}
\end{table}

The decisive comparison is the last row. Uniform INT8 produces a \emph{smaller} file than our champion
and is 39$\times$ worse on Kodak and 67$\times$ worse on CLIC. Uniform settings also vary 2.9$\times$
across datasets while mixed assignments vary at most 1.4$\times$: constraining the total rather than
each component buys robustness as well as quality.

\subsection{What QAT contributed}
Table~\ref{tab:qat} isolates the third question. QAT removes 61--98\% of the damage on assignments
chosen by hand, and \emph{nothing at all} on the assignment the solver had already optimised. The
pattern is the useful result: fine-tuning is a \emph{substitute} for allocation, not a supplement.
The size of the repair is set by what the assignment left on the table, and a mixed-integer program
that has minimised total damage leaves almost nothing.

\begin{table}[t]
\centering
\footnotesize
\setlength{\tabcolsep}{3pt}
\caption{QAT at 12{,}000 steps. Worst-setting penalty (\%) before and after; the last row used 8{,}000
steps.}
\label{tab:qat}
\begin{tabular}{lrrrr}
\toprule
Tier & MB & Before & After & Removed \\
\midrule
INT8, 16 comps    & 14.401 & 8.61 & \best{3.36} & 61.0\% \\
INT4, 2 blocks    & 13.572 & 119.18 & \best{20.53} & 82.8\% \\
INT4, 4 blocks    & 11.369 & 1771.71 & \best{41.84} & 97.6\% \\
ILP 2\% assignment & 18.389 & 1.936 & 1.936 & \best{0.0\%} \\
\bottomrule
\end{tabular}
\end{table}

\subsection{Why there is no speedup}
No configuration measured faster. Paired interleaved benchmarking across 50 pairs per cell gives
$+4.79\%$ latency for the mixed model (95\% CI $+1.27$ to $+9.80$) and $-0.39\%$ energy, whose
interval spans zero. The reason is arithmetic: a decode is 88\,ms, of which 80\,ms is entropy coding
and 8\,ms is the convolutional network. Quantization accelerates convolutions, and convolutions are
9\% of a decode and 27\% of a round trip. Simulated formats (FP8, INT4, FP4) carry no speed claim at
all; only FP16 and BF16 run real kernels here, and they save energy rather than time.

\section{Conclusion}
Granularity, not precision, is the lever: mixed precision is worth 39$\times$ over uniform precision
at comparable size. And QAT is a substitute for good allocation---it recovered 61--98\% of the damage
on hand-chosen assignments and 0\% on an optimal one. Both halves matter, and the second is only
visible because we measured the control case.

\subsection*{Limitations}
(a) Our checkpoint selector under-reports QAT: a 6-crop validation probe scored the best INT8 model as
no better than its start ($-0.08\%$) while the full evaluation shows 8.61\%$\to$3.36\%, so saved
checkpoints are probably not the best available. (b) Weights only; activations remain FP32, so W8A8 is
not exercised. (c) 4-bit is simulated, so its speed is unmodelled. (d) One entropy lane remains---the
largest unused latency lever. (e) QAT results are one seed on 100 images; the qualitative finding
holds across three tiers and two step counts, the exact percentages are single-run.

\begingroup
\small
\begin{thebibliography}{11}
\setlength{\itemsep}{0pt}
\bibitem{pulse} Jia et al. \emph{PULSE: a learned image codec}. arXiv:2609.18602, 2026.
\bibitem{jacob} Jacob et al. Quantization and training of neural networks for efficient
integer-arithmetic-only inference. \emph{CVPR}, 2018. arXiv:1712.05877.
\bibitem{bengio} Bengio, L\'eonard, Courville. Estimating or propagating gradients through stochastic
neurons for conditional computation. arXiv:1308.3432, 2013.
\bibitem{bjontegaard} G. Bj{\o}ntegaard. Calculation of average PSNR differences between RD-curves.
\emph{ITU-T VCEG-M33}, 2001.
\bibitem{pchip} Fritsch, Butland. A method for constructing local monotone piecewise cubic
interpolants. \emph{SIAM J. Sci. Comput.} 5(2):300--304, 1984. DOI 10.1137/0905021.
\bibitem{highs} Huangfu, Hall. Parallelizing the dual revised simplex method. \emph{Math. Prog.
Comp.} 10(1):119--142, 2018. DOI 10.1007/s12532-017-0130-5.
\bibitem{kodak} Eastman Kodak Company. Kodak PhotoCD image dataset, 1991.
\bibitem{clic} Toderici et al. Workshop and challenge on learned image compression. \emph{CVPR
Workshops}, 2020.
\bibitem{div2k} Agustsson, Timofte. NTIRE 2017 challenge on single image super-resolution.
\emph{CVPR Workshops}, 2017.
\bibitem{set14} Zeyde, Elad, Protter. On single image scale-up using sparse-representations.
\emph{Curves and Surfaces}, 2010.
\bibitem{code} Our code and results:
\url{https://github.com/harshith0518/CS6886-PULSE-Quantization}.
\end{thebibliography}
\endgroup

\end{document}
